

We evaluate ChatEval on two benchmarks, \textit{FairEval} and \textit{Topical-Chat} which represent the categories of open-ended question answer and dialogue response generation, respectively.


\subsection{Implementation Details}

We choose to utilize models from OpenAI's GPT family as our LLMs in ChatEval, including GPT-4 and ChatGPT (GPT-3.5-turbo) and set the temperature to 0 to ensure reproducibility. The rationale behind this selection is the exceptional performance these models offer, being among the most advanced and powerful in the world. Additionally, their accessibility and ease of use through APIs enable us to directly call and interact with the models during our research, significantly simplifying the process. In our current research, we focus on homogeneous groups of LLMs. That is, within a given multi-agent group, all LLMs belong to the same GPT family model, either all GPT-4 or all ChatGPT. We acknowledge the potential of heterogeneous groups for future research, which could provide fascinating insights into how strong models and weak models can cooperate in a multi-agent setting.

\subsection{Benchmarks}


The detailed introduction of different categories and benchmarks are listed as follows:

\textbf{Open-ended Question Answer} is a key component within the field of NLP and generative AI. It necessitates an AI system to provide comprehensive, detailed, and human-like responses to questions that don't have a predefined or fixed set of possible answers. The work of~\cite{chiang2023vicuna} encompasses a collection of 80 open-ended questions originating from a wide array of categories, including common-sense, counterfactual, coding, etc. We then take the human annotation results from~\cite{wu2023large} to conduct the experiments in this paper. For each question, they direct three annotators to evaluate the replies given by Vicuna-13B and ChatGPT through the given rules and finally derive the results by the majority votes among the annotators.


\textbf{Dialogue Response Generation} is a task involves creating a coherent and contextually appropriate response to a given input dialogue. We draw upon the \textit{Topical-Chat}~\citep{gopalakrishnan2019topical} dataset for our study. We then take the human annotation results from~\cite{mehri2020usr} where they carry out the annotations on 60 dialogue contexts with each response generated by 6 different systems. Human evaluators analyzed these responses based on \textit{natural}, \textit{coherence}, \textit{engagingness}, \textit{groundedness}, and \textit{understandable}, where we take the first four dimensions for experiments in our paper following~\cite{zhong2022towards}.


\subsection{Baselines}

We evaluate ChatEval against following methods. As the main portion of our comparison, we primarily focuses on the single-agent-based method. \textbf{Single-Agent} means that we directly query an LLM to generate the response towards the evaluation\footnote{We use the same prompt template as our multi-agent debate settings in single-agent baseline except that we ignore some slot.}. We use \textbf{Multi-Agent} to represent ChatEval where several agents discuss towards the evaluation. By default, we configure the communication strategy to one-by-one, agent numbers to 2, and discussion turns to 2 in this section and employ position calibration techniques in both single-agent and multi-agent settings. We will discuss more debate configurations in Section~\ref{sec:Analysis} for completeness. For the open-ended question answer task, we also compare our method with \textbf{FairEval}~\citep{wang2023large}. They propose various strategies to improve the evaluation performance of a LLM including Multiple Evidence Calibration (MEC) and Balanced Position Calibration (BPC). For the dialogue response generation task, we also compare our method with \textbf{G-EVAL}~\citep{liu2023gpteval}. They utilize CoT and probability-weighted summation for their method. Additionally, we include results from n-gram-based metrics, such as \textbf{ROUGE}~\citep{lin2004rouge}, \textbf{BLEU}~\citep{papineni2002bleu} and embedding-based metrics such as \textbf{BERTScore}~\citep{zhang2019bertscore}.



\subsection{Results for Open-ended question answers}\label{subsec:open_ended}

We adopt the same evaluation approach as~\cite{wang2023large} to assess the annotation results produced by different methods and annotators. Specifically, we calculate the Accuracy (Acc.), which measures the proportion of correctly classified instances out of the total instances, and the Kappa correlation coefficient (Kap.)~\citep{mchugh2012interrater} which gauges the agreement between results from models and human annotators while taking into account the possibility of agreement occurring by chance. Both metrics provide insights into the reliability and consistency of the annotations. We take the human annotation results and FairEval's~\citep{wang2023large} best results from their paper. As is shown in Table~\ref{tab:faireval_results}, different annotators can reach a relatively high agreement and perform better than any other LLM-based approach. Still, the average human annotations accuracy which is 71.7\% shows there exists a certain degree of discrepancy among different unique individuals revealing that text evaluation is absolutely an arduous task. The second part and the third part of Table~\ref{tab:faireval_results} show the results of FairEval's method and the results of our proposed method respectively.
We find that (1) ChatEval can enhance the performance of the evaluation process, achieving higher alignment with human preference compared with single-agent evaluation. Specifically, the multi-agent-based method improves the accuracy by 6.2\% for ChatGPT and 2.5\% for GPT-4; (2) ChatEval surpasses FairEval's best results within both ChatGPT and GPT-4 settings showing the effectiveness of our proposed method.


\subsection{Results for Dialogue Response Generation}
For the dialogue response generation benchmarks, we align the evaluation method with~\cite{zhong2022towards}, calculating the turn-level Spearman and Kendall-Tau correlation in correspondence with human judgments on four aspects (\textit{naturalness}, \textit{coherence}, \textit{engagingness} and \textit{groundedness}). Results can be found in Table~\ref{tab:topical_chat_results}. In the first part of Table~\ref{tab:topical_chat_results}, we demonstrate that n-gram-based metrics and embedding-based metrics perform overall poorly on all the aspects evaluated illustrating that these methods can hardly reveal human preference. In the second part of Table~\ref{tab:topical_chat_results}, we show the results from the G-eval~\citep{liu2023gpteval} paper. They first ask the LLM to generate intermediate thought and finally calculate the weighted summation of the output scores based on the probability. The results show that their method outperforms previous traditional metrics depicting the fact that the LLM-based evaluator is effective and reliable for evaluating the dialogue response generation task. While their method delivers sound results, our proposed approach raises the bar in terms of performance for GPT-4. Specifically, ChatEval improves the average Spearman and Kendall-Tau correlation by 0.096 (16.3\%) and 0.057 (10.0\%) respectively. Additionally, compared with the single-agent method, ChatEval amplifies the performance both for ChatGPT and GPT-4, showing the effectiveness of our method which is aligned with the results in Section~\ref{subsec:open_ended}.


\begin{table}[t]
\caption{Accuracy (Acc.) and Kappa correlation coefficient (Kap.) of different methods on FairEval benchmark. }
\label{tab:faireval_results}
\begin{center}
\begin{tabular}{l|l|l|l}
\hline
\textbf{Evaluator}  & \textbf{Methods}      & \textbf{Acc. (\%)} & \textbf{Kap.} \\ \hline
\multicolumn{4}{l}{\textbf{Human}}                    \\ 
Annotator1 & \multicolumn{1}{c|}{-}    & 68.8      & 0.5  \\
Annotator2 & \multicolumn{1}{c|}{-}    & 76.3      & 0.62 \\ 
Annotator3 & \multicolumn{1}{c|}{-}    & 70        & 0.5  \\ \hline
\multicolumn{4}{l}{\textbf{FairEval}}                 \\ 
ChatGPT    & MEC+BPC      & 58.7      & 0.31 \\
GPT-4      & MEC+BPC      & 62.5      & 0.37 \\ \hline
\multicolumn{4}{l}{\textbf{Ours}}                     \\ 
ChatGPT    & Single-Agent & 53.8      &   0.27   \\
ChatGPT    & Multi-Agent  & \textbf{60.0}      & \textbf{0.33}     \\
GPT-4      & Single-Agent & 61.3         & 0.36     \\
GPT-4      & Multi-Agent  & \textbf{63.8}      & \textbf{0.40}     
\end{tabular}
\end{center}
\end{table}



\begin{table}[t]
\caption{Turn-level Spearman ($\rho$) and Kendall-Tau ($\tau$) correlations of different methods on Topical-Chat benchmark, \textbf{SA} means Single-Agent and \textbf{MA} means Multi-Agent. Our ChatGPT settings should be compared to G-EVAL-3.5, and GPT-4 settings should be compared to G-EVAL-4. }
\label{tab:topical_chat_results}
\begin{center}
\begin{tabularx}{\textwidth}{p{2cm} | X | X | X | X | X | X | X | X | X | X}
\hline
\multicolumn{1}{c}{\multirow{2}{*}{Metrics}} & \multicolumn{2}{c}{Naturalness} & \multicolumn{2}{c}{Coherence} & \multicolumn{2}{c}{Engagingness} & \multicolumn{2}{c}{Groundedness} & \multicolumn{2}{c}{Average} \\
\multicolumn{1}{c}{} & $\rho$ & $\tau$ & $\rho$ & $\tau$ & $\rho$ & $\tau$ & $\rho$ & $\tau$ & $\rho$ & $\tau$\\ \hline
ROUGE-L                                      & 0.146          & 0.176          & 0.203         & 0.193         & 0.300             & 0.295          & 0.327           & 0.310           & 0.244        & 0.244       \\
BLEU-4                                       & 0.175          & 0.180           & 0.235         & 0.131         & 0.316           & 0.232          & 0.310            & 0.213          & 0.259        & 0.189        \\
BERTScore                                    & 0.209          & 0.226          & 0.233         & 0.214         & 0.335           & 0.317          & 0.317           & 0.291          & 0.274       & 0.262        \\
\hline
G-EVAL-3.5                                   & 0.539          & 0.532          & 0.544         & 0.519         & 0.691           & 0.660           & 0.567           & 0.586          & 0.585      & 0.574      \\
G-EVAL-4                                     & 0.565          & 0.549          & 0.605         & \textbf{0.594}         & 0.631           & 0.627          & 0.551           & 0.531          & 0.588        & 0.575      \\
\hline
ChatGPT(SA)                       & 0.474          & 0.421          & 0.527         & 0.482         & 0.599           & 0.549          & 0.576           & 0.558          & 0.544        & 0.503       \\
ChatGPT(MA)                        & 0.441          & 0.396          & 0.500           & 0.454         & 0.664           & 0.607          & 0.602           & 0.583          & 0.552      & 0.510         \\
\hline
GPT-4(SA)                         & 0.532          & 0.483          & 0.591         & 0.535         & 0.734           & 0.676          & \textbf{0.774}           & \textbf{0.750}           & 0.658      & 0.611        \\
GPT-4(MA)                          & \textbf{0.630}           & \textbf{0.571}          & \textbf{0.619}         & 0.561         & \textbf{0.765}           & \textbf{0.695}          & 0.722           & 0.700            & \textbf{0.684}        & \textbf{0.632}    
\end{tabularx}
\end{center}
\end{table}